\section{Dataset and Exploratory Analysis}
\label{sec:eda}

This section reports the exploratory analysis of the panel before any modelling. The aim is twofold. First, to establish the basic statistical shape of the data so that later modelling choices can be motivated rather than guessed. Second, to surface three structural findings about the panel --- multi-horizon coupling, the cumulation hypothesis, and per-horizon weight dominance --- that the midway report did not document but that materially change the modelling space. All numerical results below are produced by notebook \texttt{00\_data\_structure.ipynb} and the corresponding figures are saved as both PNG and PDF assets.

\subsection{Schema, panel scale, and train--test comparability}
\label{sec:eda:schema}

The training panel contains 5{,}337{,}414 rows, organised by 23 distinct \texttt{code} values, 180 distinct \texttt{sub\_code} values, and 5 \texttt{sub\_category} values across 4 horizons. The integer time index $\mathtt{ts\_index}$ runs over the contiguous range $[1, 3601]$. The test partition continues from $\mathtt{ts\_index} = 3602$ to 4376 with 1{,}447{,}107 rows, so test is a strict temporal suffix of train. Treating the 4-tuple $(\mathtt{code},\mathtt{sub\_code},\mathtt{sub\_category},\mathtt{horizon})$ as the natural series identifier yields 36{,}923 distinct series.

The only columns present in train but not in test are $y_{\text{target}}$ and $\mathtt{weight}$, as expected. Every category level seen in test exists in train, so there is no cold-start schema issue. The sub-code support narrows from 180 in train to 47 in test --- a coverage shift rather than a schema mismatch.

\subsection{Missing values}
\label{sec:eda:missing}

Missingness is computed on a 500{,}000-row stratified sample to keep memory in check. Of the 86 anonymised features, 48 contain at least one missing value, but the maximum missing rate across all columns is approximately 12.5\%. \Cref{fig:missingness} shows the top fifteen offenders.

\begin{figure}[h]
\centering
\includegraphics[width=0.85\linewidth]{00_missingness_top15}
\caption{Top fifteen features by missingness, on a 500{,}000-row stratified sample of train. The maximum missing rate is approximately 12.5\%, low enough that median imputation inside a fold is sufficient and there is no need to drop columns or use complex imputers.}
\label{fig:missingness}
\end{figure}

\paragraph{Reading.} Missing rates are bounded; within-fold median imputation suffices. No separate imputation strategy is developed.

\subsection{Target distribution}
\label{sec:eda:target}

The target $y_{\text{target}}$ is centred very close to zero (median $-0.0006$, IQR $\pm 0.05$), but has extremely heavy tails. The 1\% and 99\% quantiles sit at approximately $-82.8$ and $+62.9$ respectively, while the full range runs from approximately $-2202$ to $+2314$.

\begin{figure}[h]
\centering
\includegraphics[width=0.95\linewidth]{00_target_distribution}
\caption{Raw target distribution (left) and the same distribution clipped to $[-10, 10]$ (right). Most of the mass concentrates near zero, with rare extreme tails up to the order of $\pm 2200$.}
\label{fig:target_dist}
\end{figure}

\paragraph{Reading.} The shape disqualifies plain unweighted MSE as a sensible loss. Tails are real outcomes, not corruption, so we cannot simply clip them; combined with the weighting scheme below, this is the central methodological constraint.

\subsection{Weight distribution and metric concentration}
\label{sec:weight_dist}

The competition weights are extraordinarily skewed. The median weight is $1{,}699$, while the maximum is $1.39 \times 10^{13}$. \Cref{fig:weight_dist} shows the weight distribution clipped at the 99th percentile (left) and the cumulative-share Lorenz-style curve (right). The top 1\% of rows carry 64.2\% of total weight; the top 5\% carry 92.6\%; the top 10\% carry 98.3\%.

\begin{figure}[h]
\centering
\includegraphics[width=0.95\linewidth]{00_weight_distribution}
\caption{Weight distribution clipped at the 99th percentile (left) and cumulative share of total weight ranked by weight (right). Approximately 64\% of total weight is concentrated in the top 1\% of rows; 92\% in the top 5\%.}
\label{fig:weight_dist}
\end{figure}

\paragraph{Reading.} Unweighted evaluation describes a different competition. The weighted skill score (\cref{sec:skill}) is our only headline metric, and a model that accurately predicts the top one percent of rows has more leverage on the leaderboard than a uniformly mediocre one.

\subsection{Series coverage and split feasibility}
\label{sec:eda:coverage}

Not every series has enough history to support time-based validation. \Cref{fig:series_length} shows the distribution of per-series lengths. Out of 36{,}923 distinct series, only 1{,}930 (approximately 5\%) cleanly cross the train/validation cutoff at $\mathtt{ts\_index} = 2880$. Many series end before the cutoff, contributing zero validation rows.

\begin{figure}[h]
\centering
\includegraphics[width=0.85\linewidth]{00_series_length_distribution}
\caption{Distribution of per-series lengths. Most series are short, and the tail of long series carries the structural signal that classical fits can exploit.}
\label{fig:series_length}
\end{figure}

\paragraph{Reading.} Per-series classical fitting is feasible only on the eligible subset (cross the cutoff with at least 120 training rows). Global pooled models are not length-constrained --- one reason the comparison between local classical and global ML approaches is not perfectly apples-to-apples.

\subsection{Representative-series selection}
\label{sec:eda:repseries}

For every per-series figure in this project we plot the same four series. They are selected with explicit reasons rather than at random, in order to span four distinct stress regimes: long history, large weight, high variability, and near-flat behaviour. \Cref{tab:representative} lists the four chosen series. They match the selection used in the midway report, preserving visual continuity.

\begin{table}[h]
\centering
\caption{Four representative series chosen for per-series studies. Each series probes a different regime of the panel.}
\label{tab:representative}
\begin{tabular}{l l l c}
\toprule
Reason & Series key (code/sub\_code) & Horizon & Length \\
\midrule
Longest history       & X9BZ68VQ / OYJGNSQK / DPPUO5X2 & 1  & 212 \\
Highest total weight  & SJZP0OVU / OYJGNSQK / NQ58FVQM & 25 & 156 \\
Most volatile         & W4S29LF4 / KL66VIS3 / PHHHVYZI & 25 & 162 \\
Most stable           & SJZP0OVU / OYJGNSQK / NQ58FVQM & 1  & 170 \\
\bottomrule
\end{tabular}
\end{table}

\Cref{fig:rep_series_raw} shows the raw $y_{\text{target}}$ trajectory of each chosen series with the train/validation cutoff marked. The four panels look profoundly different: the most-volatile series swings between $-500$ and $+500$, while the most-stable series occupies a numerical range narrower than $0.001$.

\begin{figure}[h]
\centering
\includegraphics[width=0.95\linewidth]{00_representative_series_raw}
\caption{Raw $y_{\text{target}}$ for the four chosen representative series. The dashed vertical line marks $\mathtt{ts\_index} = 2880$, the train/validation cutoff. The four panels span the four stress regimes that any model has to handle.}
\label{fig:rep_series_raw}
\end{figure}

\paragraph{Reading.} A one-size-fits-all model will struggle: the most-stable series is essentially numerical noise, while the most-volatile series has obvious cyclical structure. Forecast performance across these four panels is one of the cleanest diagnostics for whether a model family generalises or only fits one regime.

\subsection{Stationarity diagnostics}
\label{sec:eda:stationarity}

Classical ARIMA-family forecasting assumes stationarity, or stationarity-after-differencing. We therefore run the Augmented Dickey--Fuller test on a stratified sample of 200 eligible series. The ADF test takes the null hypothesis that the series has a unit root (non-stationary); we reject at the 5\% level when the $p$-value is below $0.05$. \Cref{fig:adf} summarises the results.

\begin{figure}[h]
\centering
\includegraphics[width=0.95\linewidth]{00_adf_results}
\caption{ADF results on 200 sampled eligible series. Left: $p$-value distribution with the $0.05$ threshold marked. Centre: fraction stationary by horizon. Right: fraction stationary by code.}
\label{fig:adf}
\end{figure}

Approximately 70.5\% of sampled series reject the unit-root null at the 5\% level. The picture is more nuanced when broken down by horizon: H1 series are more often stationary than longer-horizon cumulative series, which is consistent with the multi-horizon coupling we will document in \cref{sec:cumulation}.

ADF and KPSS have opposite null hypotheses, so running both gives a more honest joint verdict: ADF takes non-stationarity as the null, while KPSS takes stationarity as the null. We classify a series as \emph{stationary} when ADF rejects and KPSS does not, as \emph{unit root} when ADF fails to reject and KPSS rejects, and as \emph{inconclusive} otherwise. \Cref{fig:adf_kpss} reports the joint distribution.

\begin{figure}[h]
\centering
\includegraphics[width=0.7\linewidth]{00_adf_kpss_joint}
\caption{Joint ADF $+$ KPSS verdict on 200 sampled series. Stationary 64.0\%, inconclusive 29.5\%, unit root 6.5\%.}
\label{fig:adf_kpss}
\end{figure}

The stationary share dominates, but the inconclusive bucket is large enough that one-size-fits-all claims about raw levels would be unsafe. We need a repair step (differencing) for the unit-root and inconclusive cases.

\Cref{fig:differencing} re-runs the ADF test on the first-differenced series. The stationary share rises from 70.5\% on raw levels to 98.5\% after differencing.

\begin{figure}[h]
\centering
\includegraphics[width=0.7\linewidth]{00_first_differencing_share}
\caption{ADF stationary share at the 5\% level: levels (left, $d=0$) versus first differences (right, $d=1$). Differencing pushes the share to nearly 100\%.}
\label{fig:differencing}
\end{figure}

\paragraph{Reading.} The diagnostic chain motivates ARIMA with $d=1$ as the default for the classical study and justifies engineered differenced features for global ML models. The 98.5\% post-difference stationarity also bears on the cumulation finding in \cref{sec:cumulation}: differencing a higher-horizon target partially inverts the rolling-sum construction back into the H1 process.

\subsection{Visual stationarity --- rolling statistics}
\label{sec:eda:rolling}

Statistical tests can disagree on short or noisy series, so a rolling-window view is a useful visual check. \Cref{fig:rolling_raw} shows a 24-step rolling mean and rolling standard deviation on the raw four representative series. The most-volatile series shows clear drift in level and bulging variance; the longest-history series drifts upward; the high-weight and most-stable series occupy near-flat ranges where the rolling mean barely moves.

\begin{figure}[h]
\centering
\includegraphics[width=0.95\linewidth]{00_rolling_raw}
\caption{Rolling mean and rolling standard deviation (window = 24) on the raw four representative series. The most-volatile series shows clear drift in level and bulging variance.}
\label{fig:rolling_raw}
\end{figure}

\Cref{fig:rolling_diff} repeats the same diagnostic on the first-differenced versions of the same series. The rolling means now hug zero and the rolling standard deviations flatten across the cutoff, visually confirming the statistical jump from the previous subsection.

\begin{figure}[h]
\centering
\includegraphics[width=0.95\linewidth]{00_rolling_differenced}
\caption{Rolling mean and rolling standard deviation (window = 24) on the first-differenced four representative series. The mean oscillates close to zero and the variance is stable, consistent with stationarity.}
\label{fig:rolling_diff}
\end{figure}

\paragraph{Reading.} Differenced series behave like the white-noise regime we want before fitting a differenced classical model.

\subsection{Lag structure --- ACF and PACF}
\label{sec:eda:acf}

The autocorrelation function (ACF) shows correlation between the series and lagged copies of itself; the partial autocorrelation function (PACF) isolates each lag's direct contribution after controlling for shorter lags. A sharp PACF cutoff suggests an autoregressive order; a sharp ACF cutoff suggests a moving-average order. Rather than relying on a single anchor series, we aggregate the absolute ACF and PACF across all 200 sampled eligible series (\cref{fig:acf_panel}). The panel-level view exposes which lags are reliably informative across the panel rather than artefacts of one trajectory.

\begin{figure}[h]
\centering
\includegraphics[width=0.95\linewidth]{00_acf_pacf_panel}
\caption{Panel-level mean $|\text{ACF}|$ (left) and mean $|\text{PACF}|$ (right) across 200 sampled eligible series. Persistent low-order lags dominate; higher lags decay quickly. There are no clean seasonal bumps.}
\label{fig:acf_panel}
\end{figure}

\paragraph{Reading.} Short-memory dynamics dominate (lags 1--3 carry most of the structure with rapid decay thereafter), and there are no clean seasonal bumps. The lag feature set for the global ML models therefore prioritises lags 1, 2, 3, 5, 10 and skips seasonal indicators.

\subsection{STL decomposition}
\label{sec:eda:stl}

Seasonal-Trend decomposition using LOESS (STL) splits a series into trend, seasonal, and residual components. We apply STL to the longest-history anchor with a period of 24 and observe the relative amplitudes of the three components. \Cref{fig:stl} shows the result.

\begin{figure}[h]
\centering
\includegraphics[width=0.85\linewidth]{00_stl_anchor}
\caption{STL decomposition of the longest-history anchor with period 24. The trend component captures slow drift; the seasonal amplitude is small relative to residual standard deviation.}
\label{fig:stl}
\end{figure}

\paragraph{Reading.} The trend captures slow drift, but the seasonal amplitude is small relative to the residual. Combined with the absence of seasonal bumps in the panel ACF, this confirms periodic structure is not the dominant signal --- modelling effort goes into short-memory dynamics, not seasonal indicators.

\subsection{Feature audit}
\label{sec:eda:feature_audit}

The 86 \texttt{feature\_*} columns are anonymised and we cannot interpret them by name. We can, however, audit them statistically. On a 200{,}000-row stratified sample we compute Pearson correlation between each feature and $y_{\text{target}}$, then rank features by absolute correlation. \Cref{fig:feature_corr} shows the top twenty.

\begin{figure}[h]
\centering
\includegraphics[width=0.85\linewidth]{00_feature_corr_top20}
\caption{Top twenty features by absolute correlation with $y_{\text{target}}$, computed on a 200{,}000-row stratified sample. Green bars are positive correlations, red bars are negative.}
\label{fig:feature_corr}
\end{figure}

\paragraph{Reading.} The single largest absolute correlation is below 0.10. No single feature predicts the target well on its own. A separate correlation check among the top twenty features further reveals strong feature--feature blocks, indicating real redundancy in the engineered set. Together these say: the signal is weak, spread across many correlated features, and best exploited by gradient-boosted trees combined with caution when reading importance plots.

\subsection{Cross-series relationships}
\label{sec:eda:cross_series}

Within the same \texttt{code}, different sub-codes might share signal. If they do, global pooled models can borrow strength across related sub-codes; per-series local models cannot. We compute lagged cross-correlations between random pairs of sub-codes within the same \texttt{(code, horizon)} group. \Cref{fig:cross_series} shows six representative pairs.

\begin{figure}[h]
\centering

\includegraphics[width=0.95\linewidth]{00_cross_series}
\caption{Lagged cross-correlation between random sub-code pairs sharing the same code and horizon. A peak away from lag 0 means one sub-code leads or lags the other.}
\label{fig:cross_series}
\end{figure}

\paragraph{Reading.} Clear lead/lag peaks within a code mean global models can learn a shared cross-sub-code component that per-series fits cannot. This is one of the strongest a-priori arguments for the global ML approach.

\subsection{Multi-horizon coupling}
\label{sec:multihorizon}

The midway report treated $(\mathtt{code},\mathtt{sub\_code},\mathtt{sub\_category},\mathtt{horizon})$ as the natural series identifier and fit each horizon independently. The exploratory analysis in this project shows that the four horizons of the same triple are not independent. For approximately 1.21 million distinct $(\mathtt{code},\mathtt{sub\_code},\mathtt{sub\_category},\mathtt{ts\_index})$ tuples, all four horizons are simultaneously present.

A second consequence of multi-horizon coupling is per-horizon target scaling. \Cref{fig:horizon_scale} shows the target standard deviation and maximum row weight by horizon.

\begin{figure}[h]
\centering
\includegraphics[width=0.95\linewidth]{00_horizon_scale_and_weight}
\caption{Per-horizon target standard deviation (left) and maximum row weight on a logarithmic scale (right). Target scale grows by a factor of approximately 4.5 from H1 to H25. Maximum weight is highest at H3.}
\label{fig:horizon_scale}
\end{figure}

The target standard deviation rises from approximately 11.7 at H1 to 52.8 at H25, a factor of $4.5$.

\paragraph{Reading.} This finding has direct modelling consequences. A single model trained on horizons pooled into one DataFrame, without horizon-aware standardisation, will be dominated by H25 magnitudes and predict at H25 scale. Predictions at H25 scale evaluated against H1 targets will overshoot by a factor of roughly 4.5, producing catastrophically large weighted RMSE on H1 rows. This is exactly the failure mode observed in earlier ML artefacts attached to the project --- the original XGBoost run reported a weighted RMSE more than two hundred times the naive baseline, consistent with a per-horizon-scale leak. The remedy is straightforward: train one model per horizon, or normalise the target by horizon before training.

\subsection{Cumulation hypothesis}
\label{sec:cumulation}

Given the multi-horizon coupling above, the natural hypothesis about how the four horizons relate to each other is that the target at horizon $k$ is a cumulative aggregate of the H1 target over the next $k$ steps:
\begin{equation}
\label{eq:cumulation}
  y^{H_k}_{t} \;\approx\; \sum_{i=0}^{k-1} y^{H_1}_{t+i}.
\end{equation}
We test \cref{eq:cumulation} on a stratified sample of 500 triples that have all four horizons and at least 60 H1 rows. For each triple we pivot the data so that horizons sit on columns, compute the rolling sum of the H1 series over the appropriate window, and report the Pearson correlation between $y^{H_k}_t$ and that rolling sum. \Cref{fig:cumulation} shows the distribution of correlations and the median ratio.

\begin{figure}[h]
\centering
\includegraphics[width=0.95\linewidth]{00_cumulation_hypothesis}
\caption{Test of the cumulation hypothesis on 500 stratified triples. Left: distribution of Pearson correlations between $y^{H_k}_t$ and the rolling sum of $y^{H_1}$ over $k$ steps, for $k \in \{3, 10, 25\}$. Right: distribution of the per-triple median ratio $y^{H_k}_t / \sum y^{H_1}$. Higher and tighter is better.}
\label{fig:cumulation}
\end{figure}

\Cref{tab:cumulation} reports the summary statistics. At $k=3$ the median Pearson correlation is 0.96 and the median ratio is 1.00, indicating essentially exact aggregation. The relation is slightly noisier at $k=10$ and $k=25$, with median ratios of 0.95 and 0.91, but still clearly dominant.

\begin{table}[h]
\centering
\caption{Cumulation hypothesis test summary on 500 triples.}
\label{tab:cumulation}
\begin{tabular}{c r r r r}
\toprule
$k$ & $n$ triples & mean corr & median corr & median ratio \\
\midrule
3   & 500 & 0.92 & 0.96 & 1.00 \\
10  & 500 & 0.86 & 0.94 & 0.95 \\
25  & 500 & 0.78 & 0.91 & 0.91 \\
\bottomrule
\end{tabular}
\end{table}

\paragraph{Reading.} Two consequences for the modelling work that follows. First, fit-H1-and-aggregate becomes a viable, possibly superior strategy to fitting per-horizon models independently: a single H1 model gives H3, H10, and H25 forecasts for free, and the aggregation is mathematically forced by the data construction. We test this in a dedicated H1-aggregation experiment in a later notebook. Second, differencing operations on $y^{H_{25}}$ partially recover the $y^{H_1}$ process, which explains why ARIMA with $d=1$ wins the most-volatile H25 series in the midway classical study: differencing inverts the cumulation. The midway interpretation that ARIMA was capturing a difficult pattern is partly correct; it was also undoing the cumulative construction.

\subsection{Per-horizon weight dominance}
\label{sec:weight_per_horizon}

The weight distribution analysis in \cref{sec:weight_dist} was panel-wide. Because $\mathtt{weight}$ and $\mathtt{horizon}$ are correlated, the rows that drive the weighted skill score may not be evenly distributed across horizons. \Cref{fig:weight_horizon} reports the share of total panel weight per horizon and the share of weight within the top 1\% of all weight mass.

\begin{figure}[h]
\centering
\includegraphics[width=0.95\linewidth]{00_weight_by_horizon}
\caption{Per-horizon weight breakdown. Left: each horizon's share of total panel weight. Right: share of the top 1\% of weight mass by horizon. H3 dominates both views.}
\label{fig:weight_horizon}
\end{figure}

The headline numbers are striking. The maximum panel weight ($1.39 \times 10^{13}$) is at H3, not H1. H3 alone holds 37.4\% of total panel weight and 43.6\% of the top-1\% weight mass; H1 holds 25.4\%, H10 holds 14.3\%, and H25 holds 16.8\%.

\paragraph{Reading.} A model that predicts H3 well on the small set of high-weight rows will tend to win the leaderboard regardless of how it performs on the other horizons. This was missed in the midway report. The practical consequence is that we report H3 weighted skill as a separate headline number alongside the panel-wide score, and we track H3 performance carefully when comparing models.

\subsection{EDA-driven modelling implications}
\label{sec:eda:implications}

The EDA findings above constrain every modelling chapter that follows. The classical baselines (notebook 01) honour the chronological cutoff and use the four representative series. The global LightGBM models (notebook 02) train per horizon with sample weights, motivated by the per-horizon target scaling and H3 weight dominance. The H1-aggregation experiment (notebook 03) operationalises the cumulation hypothesis. The probabilistic forecasting work (notebook 04) addresses the heavy tails and weak per-feature signal directly through quantile regression and conformal calibration. The ensembles (notebook 05) fit weights on a meta slice that respects chronology. The final comparison (notebook 06) reports both panel-wide and H3-specific weighted skill scores with bootstrap intervals.
